\chapter{Archimedean Tate Kernel Positivity: the Integral Energy Sub-Class}
\label{v4-arith:w7-e:chapter}

\emph{Chapter~\ref{v4-arith:w6-a:ATP-direct-obstruction} reduced A-TP to a
bounded-interval positivity inequality on the core
\([-t_{\ast},t_{\ast}]\), with \(t_{\ast}\) the unique positive zero of
the archimedean digamma kernel
\(g(t)=\mathrm{Re}\,\psi(\tfrac{1}{4}+\tfrac{it}{2})+\log\pi\). The bounded-interval reduction
proved \(W_{\infty}\ge0\) on the high-frequency dominance class
\(\mathrm{PW}^{\mathrm{hol},\mathrm{HFD}}_{F}\) defined through a
supremum condition on \(\widehat{f}\) over the core. A transverse
integral-energy sub-class isolates the
far-tail mass which the bounded-interval reduction supremum condition does not measure.
The chapter proves a sharp lower bound
\(g(t)\geq g_{L}(t):=\tfrac{1}{2}\log((1+4t^{2})/16)+\log\pi-2/(1+4t^{2})\)
for \(|t|\geq t_{L}^{\ast}\approx 0.8946\), derived from the certified
Stirling-4 remainder; an integral-energy sub-class
\(\mathrm{PW}^{\mathrm{hol},\mathrm{IE}}_{F}\)
defined by \(M\cdot\alpha(f)\leq g_{L}(T_{1})\cdot\beta(f,T_{1})\)
for some \(T_{1}\geq t_{L}^{\ast}\), where \(\alpha(f)\) and \(\beta(f,T_{1})\)
are the core and far-tail \(L^{2}\)-energy fractions of \(\widehat{f}\);
and an unconditional \(W_{\infty}\)-positivity theorem on this class,
together with an
explicit characterization of the
complementary sub-class
\(\mathrm{PW}^{\mathrm{hol}}_{F}\setminus\bigcup_{T_{1}}\mathrm{PW}^{\mathrm{hol},\mathrm{IE}}_{F}(T_{1})\)
by the failure of every logarithmic far-tail payment. No inclusion of
HFD in IE is used.}

\section{Kernel and Integral Energy}
\label{v4-arith:w7-e:context}

Chapter~\ref{v4-arith:w6-a:ATP-direct-obstruction} established: the
archimedean Weil contribution decomposes as
\(W_{\infty}(f*\widetilde{f})=(B[f]-A[f])/(2\pi)\) with
\(A[f]=\int_{|t|<t_{\ast}}|\widehat{f}(t)|^{2}(-g(t))\,dt\geq 0\) and
\(B[f]=\int_{|t|>t_{\ast}}|\widehat{f}(t)|^{2}g(t)\,dt\geq 0\). The
inequality \(A[f]\leq B[f]\) is exactly
\(W_{\infty}(f*\widetilde f)\ge0\); full A-TP still requires the
residual \(R[f]\) of \Cref{v4-arith:w6-a:eq:R-definition}. The bounded-interval reduction
proves the
inequality \(A[f]\leq B[f]\) on
\(\mathrm{PW}^{\mathrm{hol},\mathrm{HFD}}_{F}\) where
\(\sup_{|t|<t_{\ast}}|\widehat{f}(t)|^{2}\cdot 2t_{\ast}M\leq B[f]\).
This condition is restrictive: a sharp spike in \(\widehat{f}\) on the
core kills the sup-bound even when the integral
\(\int_{|t|<t_{\ast}}|\widehat{f}|^{2}\) is small.

\(W_{\infty}\ge0\) holds on an additional sub-class under integral
control of the core and a lower bound for \(g\) on the far tail. The
resulting class is a transverse sufficient criterion, not an ordered
replacement for HFD.

\subsection{Domain}
\label{v4-arith:w7-e:domain}

The objects are:
\begin{itemize}
\item explicit Stirling-type lower bound for \(g\) on
\(|t|\geq t_{L}^{\ast}\), valid unconditionally by the Stirling-4
remainder bound;
\item integral-energy sub-class defined through Plancherel masses
rather than supremum values;
\item unconditional \(W_{\infty}\)-positivity on this sub-class;
\item characterization of the complementary residual by failure of
all logarithmic far-tail payments, and a separate two-scale prolate
sufficient criterion.
\end{itemize}

A-TP on the full Paley--Wiener class is RH-equivalent by
\Cref{v4-arith:w5-a:thm:ATP-equals-weil-positivity}. The integral-energy
sub-class does not capture all of \(\mathrm{PW}^{\mathrm{hol}}_{F}\).

\subsection{Independent Sources}
\label{v4-arith:w7-e:sources}

Four classical sources enter the analysis.

\begin{description}
\item[Stirling 1730.] Stirling, J. \emph{Methodus Differentialis}.
London 1730, Prop.~28. Asymptotic expansion of \(\log\Gamma(z)\) and
its logarithmic derivative \(\psi(z)\), with alternating Bernoulli
tail.
\item[Bernoulli numbers (Euler 1755).] Euler, L. \emph{Institutiones
Calculi Differentialis}. St.\ Petersburg 1755, II.\ ch.~vi. Signs
and magnitudes of \(B_{2k}\); alternation property.
\item[Plancherel 1910.] Plancherel, M. \emph{Contribution \`a l'\'etude
de la repr\'esentation d'une fonction arbitraire par des int\'egrales
d\'efinies}. \emph{Rend.~Circ.~Mat.~Palermo} 30 (1910), 289--335.
Parseval identity for the Fourier transform on \(L^{2}(\mathbb{R})\).
\item[Slepian--Pollak 1961.] Slepian, D.; Pollak, H.O. \emph{Prolate
spheroidal wave functions, Fourier analysis, and uncertainty, I}.
\emph{Bell Syst.~Tech.~J.} 40 (1961), 43--63. Prolate eigenvalue
structure of the bandpass concentration operator on
\(L^{2}([-T,T])\to L^{2}([-\Omega,\Omega])\).
\end{description}

Stirling 1730, Euler 1755, Plancherel 1910,
and Slepian--Pollak 1961 do not enter Weil 1952, Li 1997, the
Rankin--Selberg integral, or the Hurwitz derivation used in
Chapter~\ref{v4-arith:w6-a:ATP-direct-obstruction}. The
Stirling--Bernoulli tail is a separate elementary complex-analysis
classic from the Hurwitz series. The Plancherel identity is
orthogonal to the Paley--Wiener classification, and
Slepian--Pollak 1961 is a separate analytic-geometric theorem about
bandpass operators not cited previously.

\section{A Sharp Stirling Lower Bound for the Archimedean Kernel}
\label{v4-arith:w7-e:stirling-lower-bound}

\subsection{Stirling Series for the Digamma and Sign Control}
\label{v4-arith:w7-e:stirling}

\begin{lemma}[Stirling series for \(\psi\) with Bernoulli tail]
\label{v4-arith:w7-e:lem:stirling-series}
\ClaimStatusProvedHere. For \(\mathrm{Re}\,z>0\) and every \(N\geq 0\),
\begin{equation}
\label{v4-arith:w7-e:eq:stirling-series}
\psi(z)\;=\;\log z-\frac{1}{2z}-\sum_{k=1}^{N}\frac{B_{2k}}{2k\,z^{2k}}+R_{N}(z),
\end{equation}
where \(B_{2k}\) are the even Bernoulli numbers (\(B_{2}=1/6\),
\(B_{4}=-1/30\), \(B_{6}=1/42\), alternating signs) and the remainder
\(R_{N}(z)\) satisfies
\begin{equation}
\label{v4-arith:w7-e:eq:remainder-bound}
|R_{N}(z)|\;\leq\;\frac{|B_{2N+2}|}{(2N+2)\,|z|^{2N+2}}\cdot\frac{1}{\cos^{2N+2}(\arg z /2)},
\quad|\arg z|<\pi.
\end{equation}
\end{lemma}

\begin{proof}
Whittaker--Watson 1927 \S12.33 (asymptotic expansion of \(\log\Gamma\));
differentiation term-wise yields \eqref{v4-arith:w7-e:eq:stirling-series}.
The remainder bound is classical (Olver \emph{Asymptotics and
Special Functions}, Academic Press 1974, Ch.~8, Thm.~8.1) applied to
\(\psi\) via the integral representation
\(\psi(z)=\log z-1/(2z)-2\int_{0}^{\infty}t\bigl((t^{2}+z^{2})(e^{2\pi t}-1)\bigr)^{-1}\,dt\).
\end{proof}

\begin{lemma}[real-part Stirling at \(z=1/4+it/2\)]
\label{v4-arith:w7-e:lem:re-psi}
\ClaimStatusProvedHere. For \(t\in\mathbb{R}\) set
\(z(t):=1/4+it/2\). Then
\begin{align}
\label{v4-arith:w7-e:eq:re-log-z}
\mathrm{Re}\,\log z(t)&\;=\;\tfrac{1}{2}\log(1/16+t^{2}/4)\;=\;\tfrac{1}{2}\log\bigl((1+4t^{2})/16\bigr),\\
\label{v4-arith:w7-e:eq:re-reciprocal-z}
\mathrm{Re}\,\frac{-1}{2z(t)}&\;=\;-\frac{1}{8|z(t)|^{2}}\;=\;-\frac{2}{1+4t^{2}}.
\end{align}
The one-step truncation of \eqref{v4-arith:w7-e:eq:stirling-series}
(with \(N=0\), no Bernoulli terms) gives
\begin{equation}
\label{v4-arith:w7-e:eq:stirling-1}
\mathrm{Re}\,\psi(z(t))\;=\;\tfrac{1}{2}\log\bigl((1+4t^{2})/16\bigr)-\frac{2}{1+4t^{2}}+\mathrm{Re}\,R_{0}(z(t)).
\end{equation}
\end{lemma}

\begin{proof}
\(|z(t)|^{2}=(1/4)^{2}+(t/2)^{2}=1/16+t^{2}/4\). Hence \(\log|z|=
(1/2)\log(1/16+t^{2}/4)\). For \(1/z\): \(\mathrm{Re}(1/z)=\mathrm{Re}\,z/|z|^{2}=(1/4)/|z|^{2}\),
so \(\mathrm{Re}(-1/(2z))=-(1/8)/|z|^{2}=-2/(1+4t^{2})\). The identity
follows by direct substitution in \eqref{v4-arith:w7-e:eq:stirling-series}.
\end{proof}

\subsection{The Lower Bound \(g_{L}\)}
\label{v4-arith:w7-e:g-lower-bound}

\begin{theorem}[Stirling lower bound for \(g\)]
\label{v4-arith:w7-e:thm:g-lower-bound}
\ClaimStatusProvedHere. Define
\begin{equation}
\label{v4-arith:w7-e:eq:g-L-def}
g_{L}(t)\;:=\;\tfrac{1}{2}\log\bigl((1+4t^{2})/16\bigr)+\log\pi-\frac{2}{1+4t^{2}}.
\end{equation}
Then \(g(t)\geq g_{L}(t)\) for every \(|t|\geq t_{L}^{\ast}\), where
\(t_{L}^{\ast}\approx 0.8946\) is the unique positive zero of \(g_{L}\).
\end{theorem}

\begin{proof}
Use the certified Stirling-4 lower bound of
\Cref{v4-arith:w7-a:def:g-lower} and
\Cref{v4-arith:w7-a:prop:g-stirling-4}. With
\(u=1+4t^{2}\),
\[
g^{-}_{\mathrm{St4}}(t)
=g_{L}(t)-\frac{4}{3}\frac{2-u}{u^{2}}-\frac{32}{15u^{2}}
=g_{L}(t)+\frac{20u-72}{15u^{2}}.
\]
At \(|t|\geq t_{L}^{\ast}\) one has \(u\geq 1+4(t_{L}^{\ast})^{2}>18/5\).
Thus \(g^{-}_{\mathrm{St4}}(t)\geq g_{L}(t)\), and
\(\Cref{v4-arith:w7-a:def:g-lower}\) gives \(g(t)\geq g_{L}(t)\) on this
tail. No assertion is made on the core, where \(g_{L}\) is not a lower
bound.
\end{proof}

\begin{remark}[quantitative gap]
\label{v4-arith:w7-e:rem:quantitative-gap}
Numerically, \(g(t)-g_{L}(t)\approx 0.10\) for \(t\in[1,2]\) and
\(g(t)-g_{L}(t)\approx 0.013\) at \(t=5\), decaying as \(O(t^{-2})\) by
the next Bernoulli term. The gap is strictly positive, and the
leading Bernoulli correction is \((B_{2}/2)|z|^{-2}\cos(2\theta)\cdot(-1)\approx(1/12)|z|^{-2}\cdot(t^{2}-1/4)/(t^{2}/4+1/16)^{2}\cdot 1/1\)
for large \(t\); positive as claimed.
\end{remark}

\begin{corollary}[\(g_{L}\) vanishes at an explicit point]
\label{v4-arith:w7-e:cor:g-L-zero}
\ClaimStatusProvedHere. There is a unique \(t_{L}^{\ast}>0\) with
\(g_{L}(t_{L}^{\ast})=0\). Numerically \(t_{L}^{\ast}\approx 0.8946\);
in particular \(t_{L}^{\ast}>t_{\ast}\approx 0.8507\). For
\(|t|\geq t_{L}^{\ast}\), \(g_{L}(t)\geq 0\), and \(g(t)\geq g_{L}(t)\geq 0\).
\end{corollary}

\begin{proof}
\(g_{L}(0)=\log(1/16)/2+\log\pi-2=-\log 4+\log\pi-2\approx-2.24<0\).
As \(|t|\to\infty\), \(g_{L}(t)\to\log(|t|\pi/2)-0\to\infty\). Continuity
and monotonicity (derivative direct) supply existence and uniqueness
of \(t_{L}^{\ast}\). Numerical bisection: \(g_{L}(0.89)\approx-0.008\),
\(g_{L}(0.90)\approx 0.009\); \(t_{L}^{\ast}\approx 0.8946\).
Strict inequality \(t_{L}^{\ast}>t_{\ast}\) follows from \(g(t_{\ast})=0\)
and \(g(t)>g_{L}(t)\).
\end{proof}

\subsection{Sign-Domain Boundary}
\label{v4-arith:w7-e:sign-domain}

\begin{center}
\small
\begin{tabular}{c|ccc}
Interval & \(g(t)\) & \(g_{L}(t)\) & \(g(t)-g_{L}(t)\) \\\hline
\(|t|<t_{\ast}\) & \(<0\) & \(<0\) & sign varies \\
\(t_{\ast}\leq|t|<t_{L}^{\ast}\) & \(\geq 0\) & \(<0\) & \(>0\) \\
\(|t|\geq t_{L}^{\ast}\) & \(>0\) & \(\geq 0\) & \(\geq 0\) \\
\end{tabular}
\end{center}

On the intermediate band \([t_{\ast},t_{L}^{\ast}]\), \(g\geq 0\) but
\(g_{L}<0\); the Stirling lower bound is unusable there and we must use
the direct \(g\geq 0\) bound from monotonicity
(\Cref{v4-arith:w6-a:thm:t-star}). On \(|t|\geq t_{L}^{\ast}\), both
bounds are positive.

\section{The Integral-Energy Sub-Class}
\label{v4-arith:w7-e:integral-energy}

\subsection{Energy Fractions}
\label{v4-arith:w7-e:energy-fractions}

\begin{definition}[energy fractions]
\label{v4-arith:w7-e:def:energy-fractions}
\ClaimStatusDefinitional. Fix \(T_{1}\geq t_{L}^{\ast}\). For
\(f\in\mathrm{PW}(\mathbb{R})\) with \(\widehat{f}\not\equiv 0\), define
\begin{align}
\label{v4-arith:w7-e:eq:alpha-core}
\alpha(f)&\;:=\;\frac{\int_{|t|<t_{\ast}}|\widehat{f}(t)|^{2}\,dt}{\int_{\mathbb{R}}|\widehat{f}(t)|^{2}\,dt},\\
\label{v4-arith:w7-e:eq:beta-far}
\beta(f,T_{1})&\;:=\;\frac{\int_{|t|>T_{1}}|\widehat{f}(t)|^{2}\,dt}{\int_{\mathbb{R}}|\widehat{f}(t)|^{2}\,dt},\\
\label{v4-arith:w7-e:eq:gamma-mid}
\gamma(f,T_{1})&\;:=\;1-\alpha(f)-\beta(f,T_{1}).
\end{align}
All three are in \([0,1]\) and sum to \(1\). \(\alpha\) is the fraction of
spectral mass on the core; \(\beta\) on the far tail; \(\gamma\) on the
intermediate band \(t_{\ast}\leq|t|\leq T_{1}\).
\end{definition}

\subsection{Integral-Energy Sub-Class \(\mathrm{PW}^{\mathrm{hol},\mathrm{IE}}_{F}\)}
\label{v4-arith:w7-e:IE-definition}

\begin{definition}[integral-energy sub-class]
\label{v4-arith:w7-e:def:IE-subclass}
\ClaimStatusDefinitional. Fix \(T_{1}\geq t_{L}^{\ast}\). Define
\begin{equation}
\label{v4-arith:w7-e:eq:IE-condition}
\mathrm{PW}^{\mathrm{hol},\mathrm{IE}}_{F}(T_{1})\;:=\;\Bigl\{f\in\mathrm{PW}^{\mathrm{hol}}_{F}\;:\;
M\cdot\alpha(f)\;\leq\;g_{L}(T_{1})\cdot\beta(f,T_{1})\Bigr\},
\end{equation}
with \(M=|g(0)|\approx 3.083\) and \(g_{L}\) as in
\eqref{v4-arith:w7-e:eq:g-L-def}. The sub-class is parametrized by
the choice of \(T_{1}\geq t_{L}^{\ast}\); different \(T_{1}\) capture
different portions of \(\mathrm{PW}^{\mathrm{hol}}_{F}\).
\end{definition}

\begin{remark}[monotonicity in \(T_{1}\)]
\label{v4-arith:w7-e:rem:monotonicity-T1}
The condition \eqref{v4-arith:w7-e:eq:IE-condition} is not monotone
in \(T_{1}\): increasing \(T_{1}\) enlarges \(g_{L}(T_{1})\) (the left-hand bound
tightens) but shrinks \(\beta(f,T_{1})\) (less far-tail
mass) and enlarges \(\gamma(f,T_{1})\) (more middle-band mass, which
does not contribute to the condition). The optimal \(T_{1}\) depends
on the spectral profile of \(f\).
\end{remark}

\subsection{Unconditional \(W_{\infty}\)-Positivity on \(\mathrm{PW}^{\mathrm{hol},\mathrm{IE}}_{F}\)}
\label{v4-arith:w7-e:IE-closure}

\begin{theorem}[integral-energy \(W_{\infty}\)-positivity]
\label{v4-arith:w7-e:thm:IE-ATP}
\ClaimStatusProvedHere. For every \(T_{1}\geq t_{L}^{\ast}\) and every
\(f\in\mathrm{PW}^{\mathrm{hol},\mathrm{IE}}_{F}(T_{1})\),
\begin{equation}
\label{v4-arith:w7-e:eq:IE-closure}
W_{\infty}(f*\widetilde{f})\;\geq\;0.
\end{equation}
Full A-TP follows on the sublocus where
\(R[f]\geq W_{\infty}(f*\widetilde f)\).
\end{theorem}

\begin{proof}
By \Cref{v4-arith:w6-a:lem:decomposition},
\(W_{\infty}(f*\widetilde{f})\geq 0\iff B[f]\geq A[f]\). We upper-bound
\(A[f]\) and lower-bound \(B[f]\).

\emph{Upper bound for \(A[f]\).} On the core \(|t|<t_{\ast}\), \(|g(t)|\leq M\)
by \Cref{v4-arith:w6-a:lem:g-monotone}: \(g\) is even and monotone in
\(|t|\) with \(g(0)=-M\), so \(|g(t)|\leq M\). Hence
\[
A[f]\;\leq\;M\int_{|t|<t_{\ast}}|\widehat{f}(t)|^{2}\,dt\;=\;M\cdot\alpha(f)\cdot\|\widehat{f}\|_{L^{2}}^{2}.
\]

\emph{Lower bound for \(B[f]\).} Split
\(B[f]=\int_{|t|\geq t_{\ast}}|\widehat{f}|^{2}g\,dt\) into the middle
band \([t_{\ast},T_{1}]\) and far band \([T_{1},\infty)\) (and symmetric
negative counterparts). On the middle band, \(g\geq 0\)
(\Cref{v4-arith:w6-a:thm:t-star}). On the far band,
\(g\geq g_{L}\geq g_{L}(T_{1})\geq 0\) (by monotonicity of \(g_{L}\) in
\(|t|\) for \(|t|\geq t_{L}^{\ast}\leq T_{1}\)). Hence
\[
B[f]\;\geq\;\int_{|t|>T_{1}}|\widehat{f}|^{2}g(t)\,dt\;\geq\;g_{L}(T_{1})\int_{|t|>T_{1}}|\widehat{f}|^{2}\,dt\;=\;g_{L}(T_{1})\cdot\beta(f,T_{1})\cdot\|\widehat{f}\|_{L^{2}}^{2}.
\]

Combining: \(A[f]\leq M\alpha\|\widehat{f}\|^{2}\) and
\(B[f]\geq g_{L}(T_{1})\beta\|\widehat{f}\|^{2}\). The IE condition
\eqref{v4-arith:w7-e:eq:IE-condition}
\(M\alpha\leq g_{L}(T_{1})\beta\) gives \(A[f]\leq B[f]\), hence
\(W_{\infty}(f*\widetilde{f})\geq 0\). The final assertion follows
from \Cref{v4-arith:w6-a:thm:bounded-reduction} on the sublocus
\(R[f]\geq W_{\infty}(f*\widetilde f)\).
\end{proof}

\subsection{Strict Enlargement over the bounded-interval reduction HFD}
\label{v4-arith:w7-e:strict-enlargement}

\begin{proposition}[IE is a transverse extension of HFD]
\label{v4-arith:w7-e:prop:strict-inclusion}
\ClaimStatusProvedHere. The class
\[
\mathrm{PW}^{\mathrm{hol},\mathrm{HFD}}_{F}\cup
\bigcup_{T_{1}\geq t_{L}^{\ast}}\mathrm{PW}^{\mathrm{hol},\mathrm{IE}}_{F}(T_{1})
\]
strictly contains \(\mathrm{PW}^{\mathrm{hol},\mathrm{HFD}}_{F}\). The
inclusion
\[
\mathrm{PW}^{\mathrm{hol},\mathrm{HFD}}_{F}\subset
\bigcup_{T_{1}\geq t_{L}^{\ast}}\mathrm{PW}^{\mathrm{hol},\mathrm{IE}}_{F}(T_{1})
\]
is not asserted.
\end{proposition}

\begin{proof}
The HFD condition controls the \(L^{\infty}\)-height of
\(\widehat f\) on the core; the IE condition controls only its
\(L^{2}\)-mass there, paid for by a fixed far-tail mass. Choose a
smooth compactly supported even packet \(\varphi\), dilate it so that
\(\widehat\varphi\) is concentrated in a short interval in
\((-t_{\ast},t_{\ast})\), and add a small even modulation
\(\cos(Ru)\varphi(u)\) with \(R>T_{1}\). For large dilation the core
packet has arbitrarily small \(L^{2}\)-mass and arbitrarily large
pointwise height relative to that mass. Choosing the modulation
coefficient so that
\[
M\alpha(f)<g_{L}(T_{1})\beta(f,T_{1})
\quad\text{but}\quad
2t_{\ast}M\sup_{|t|<t_{\ast}}|\widehat f(t)|^{2}>B[f]
\]
gives \(f\in\mathrm{PW}^{\mathrm{hol},\mathrm{IE}}_{F}(T_{1})\) and
\(f\notin\mathrm{PW}^{\mathrm{hol},\mathrm{HFD}}_{F}\). The construction is
open under small \(C^{\infty}_{c}\)-perturbation, so the containment in
the displayed union is strict. The reverse containment is not used.
\end{proof}

\begin{remark}[which profile IE captures and HFD misses]
\label{v4-arith:w7-e:rem:IE-advantage}
The HFD class is constrained by pointwise control of \(\widehat{f}\)
over the core; the IE class is constrained by integrated \(L^{2}\)
control and by a far-tail payment. A narrow core peak with distributed
tail may lie in IE and fail HFD. A tail concentrated near
\(t_{\ast}\) may lie in HFD and fail IE. The two conditions are
transverse sufficient criteria.
\end{remark}

\section{Complement of \(\mathrm{PW}^{\mathrm{hol},\mathrm{IE}}_{F}\)}
\label{v4-arith:w7-e:residual}

\subsection{Complement Structure}
\label{v4-arith:w7-e:complement}

\begin{theorem}[IE residual as concentration inequality]
\label{v4-arith:w7-e:thm:IE-residual}
\ClaimStatusProvedHere. \(f\in\mathrm{PW}^{\mathrm{hol}}_{F}\setminus\bigcup_{T_{1}}\mathrm{PW}^{\mathrm{hol},\mathrm{IE}}_{F}(T_{1})\)
if and only if for every \(T_{1}\geq t_{L}^{\ast}\),
\begin{equation}
\label{v4-arith:w7-e:eq:residual-concentration}
g_{L}(T_{1})\beta(f,T_{1})\;<\;M\alpha(f).
\end{equation}
In particular,
\[
\sup_{T_{1}\geq t_{L}^{\ast}}g_{L}(T_{1})\beta(f,T_{1})\leq M\alpha(f).
\]
Thus the residual is not the complement of high frequency; it is the
locus where no logarithmic far-tail threshold pays for the core.
\end{theorem}

\begin{proof}
\eqref{v4-arith:w7-e:eq:residual-concentration} is
\eqref{v4-arith:w7-e:eq:IE-condition} reversed. For residual \(f\) the
reversed inequality holds at every \(T_{1}\); conversely the reversed
inequality at every \(T_{1}\) excludes membership in every IE class.
Taking the supremum gives the displayed necessary bound. No decay statement
stronger than this weighted-tail bound follows without an independent
regularity estimate on \(f\).
\end{proof}

\begin{remark}[which test functions are in the residual]
\label{v4-arith:w7-e:rem:residual-characterization}
Residual test functions have too little logarithmically weighted
far-tail mass relative to their core mass. Low-frequency approximate
identities give model examples after the weighted tail is made
uniformly small. Individual Li approximants require a separate
spectral comparison; no membership assertion for them is made here.
\end{remark}

\subsection{Prolate-Spheroidal Characterization}
\label{v4-arith:w7-e:prolate}

\begin{theorem}[two-scale prolate sufficient criterion]
\label{v4-arith:w7-e:thm:prolate-characterization}
\ClaimStatusProvedHere. Fix Paley--Wiener exponential type \(T\) for
\(f\), i.e., \(\mathrm{supp}(f)\subset[-T,T]\). Let
\(\{\lambda_{k}(T,\Omega)\}_{k=0}^{\infty}\) denote the
Slepian--Pollak prolate eigenvalues of the integral operator
\(K_{T,\Omega}\colon L^{2}([-T,T])\to L^{2}([-T,T])\)
\begin{equation}
\label{v4-arith:w7-e:eq:prolate-operator}
(K_{T,\Omega}f)(u)\;:=\;\int_{-T}^{T}\frac{\sin(\Omega(u-v))}{\pi(u-v)}\,f(v)\,dv,
\end{equation}
sorted \(\lambda_{0}\geq\lambda_{1}\geq\ldots>0\). Then for
\(f\in\mathrm{PW}(\mathbb{R})\) with support in \([-T,T]\),
\begin{equation}
\label{v4-arith:w7-e:eq:prolate-alpha-bound}
\alpha(f)\;\leq\;\lambda_{0}(T,t_{\ast}).
\end{equation}
For every \(T_{1}\geq t_{L}^{\ast}\),
\begin{equation}
\label{v4-arith:w7-e:eq:prolate-beta-bound}
\beta(f,T_{1})\;\geq\;1-\lambda_{0}(T,T_{1}).
\end{equation}
Consequently all \(f\in\mathrm{PW}^{\mathrm{hol}}_{F}\) of type at most
\(T\) lie in \(\mathrm{PW}^{\mathrm{hol},\mathrm{IE}}_{F}(T_{1})\) if
\begin{equation}
\label{v4-arith:w7-e:eq:two-scale-prolate-condition}
M\lambda_{0}(T,t_{\ast})\;\leq\;
g_{L}(T_{1})\bigl(1-\lambda_{0}(T,T_{1})\bigr).
\end{equation}
\end{theorem}

\begin{proof}
The operator \(K_{T,\Omega}\) is the time-limited frequency projection
onto \([-\Omega,\Omega]\). Hence
\[
\frac{\int_{-\Omega}^{\Omega}|\widehat f(t)|^{2}\,dt}
{\int_{\mathbb R}|\widehat f(t)|^{2}\,dt}
=\frac{\langle K_{T,\Omega}f,f\rangle}{\langle f,f\rangle}
\leq \lambda_{0}(T,\Omega).
\]
Taking \(\Omega=t_{\ast}\) gives
\eqref{v4-arith:w7-e:eq:prolate-alpha-bound}. Taking
\(\Omega=T_{1}\) gives the mass bound inside \([-T_{1},T_{1}]\), hence
\(\beta(f,T_{1})\geq1-\lambda_{0}(T,T_{1})\). If
\eqref{v4-arith:w7-e:eq:two-scale-prolate-condition} holds, then
\[
M\alpha(f)\leq M\lambda_{0}(T,t_{\ast})
\leq g_{L}(T_{1})(1-\lambda_{0}(T,T_{1}))
\leq g_{L}(T_{1})\beta(f,T_{1}),
\]
so \(f\in\mathrm{PW}^{\mathrm{hol},\mathrm{IE}}_{F}(T_{1})\).
\end{proof}

\begin{remark}[quantitative prolate bound]
\label{v4-arith:w7-e:rem:prolate-quantitative}
For the Slepian parameter \(c:=T\cdot t_{\ast}\), the top prolate
eigenvalue satisfies \(\lambda_{0}(T,t_{\ast})\approx 2c/\pi\) for
small \(c\) and \(\lambda_{0}\to 1\) as \(c\to\infty\) (Slepian--Pollak
1961 \S IV). For \(T=1\), \(c\approx 0.85\), \(\lambda_{0}\approx 0.54\);
\(T=0.5\), \(c\approx 0.43\), \(\lambda_{0}\approx 0.27\). These
single-scale numbers bound only the core mass. They do not imply IE
membership unless the second prolate number \(\lambda_{0}(T,T_{1})\)
also gives enough far-tail mass in
\eqref{v4-arith:w7-e:eq:two-scale-prolate-condition}.
\end{remark}

\subsection{Absolute Uniform Closure Sub-Class}
\label{v4-arith:w7-e:uniform-sub-class}

A uniform closure from prolate theory requires a core scale and a
far-tail scale.

\begin{corollary}[uniform \(W_{\infty}\)-positivity under a two-scale prolate inequality]
\label{v4-arith:w7-e:cor:uniform-narrow-PW}
\ClaimStatusProvedHere. Fix \(T_{0}>0\) and \(T_{1}\geq t_{L}^{\ast}\).
If
\begin{equation}
\label{v4-arith:w7-e:eq:uniform-narrow-condition}
M\lambda_{0}(T_{0},t_{\ast})\;\leq\;
g_{L}(T_{1})\bigl(1-\lambda_{0}(T_{0},T_{1})\bigr),
\end{equation}
then \(W_{\infty}(f*\widetilde{f})\geq0\) for every
\(f\in\mathrm{PW}^{\mathrm{hol}}_{F}\) of exponential type \(T\leq T_{0}\).
\end{corollary}

\begin{proof}
The top prolate eigenvalue is monotone in both the time window and the
frequency window. Hence for \(T\leq T_{0}\),
\[
\alpha(f)\leq\lambda_{0}(T_{0},t_{\ast}),\qquad
\beta(f,T_{1})\geq1-\lambda_{0}(T_{0},T_{1}).
\]
Hypothesis \eqref{v4-arith:w7-e:eq:uniform-narrow-condition} gives
\(M\alpha(f)\leq g_{L}(T_{1})\beta(f,T_{1})\). Thus
\(f\in\mathrm{PW}^{\mathrm{hol},\mathrm{IE}}_{F}(T_{1})\), and
\Cref{v4-arith:w7-e:thm:IE-ATP} gives the claim.
\end{proof}

\begin{remark}[significance]
\label{v4-arith:w7-e:rem:narrow-PW-significance}
\Cref{v4-arith:w7-e:cor:uniform-narrow-PW} is a uniform
\(W_{\infty}\)-positivity result only after the two prolate eigenvalues in
\eqref{v4-arith:w7-e:eq:uniform-narrow-condition} are verified. The
single estimate \(\alpha(f)\leq\lambda_{0}(T_{0},t_{\ast})\) controls
the core but says nothing about how much mass has crossed the far-tail
threshold \(T_{1}\).
\end{remark}

\section{Quantification of the Archimedean Positivity Extension}
\label{v4-arith:w7-e:quantification}

\begin{theorem}[\(W_{\infty}\)-positivity extent]
\label{v4-arith:w7-e:thm:extent}
\ClaimStatusProvedHere. Let \(\mathcal{S}\subset\mathrm{PW}^{\mathrm{hol}}_{F}\)
denote the set of real even Paley--Wiener test functions (satisfying
H1--H3 of \Cref{v4-arith:w5-a:lem:upsilon-cuspidal}). Then
\begin{equation}
\label{v4-arith:w7-e:eq:extent-bound}
\mathrm{PW}^{\mathrm{hol},\mathrm{HFD}}_{F}\;\subsetneq\;
\mathrm{PW}^{\mathrm{hol},\mathrm{HFD}}_{F}\cup
\bigcup_{T_{1}}\mathrm{PW}^{\mathrm{hol},\mathrm{IE}}_{F}(T_{1})
\;\subsetneq\;\mathcal{S}.
\end{equation}
The first inclusion is strict by
\Cref{v4-arith:w7-e:prop:strict-inclusion}. The second is strict:
there exist smooth compactly supported even test functions whose
Fourier mass is concentrated in the core and whose logarithmically
weighted far tail is arbitrarily small.
\end{theorem}

\begin{proof}
First inclusion: \Cref{v4-arith:w7-e:prop:strict-inclusion}.
For the second inclusion, take an even non-zero
\(\varphi\in C^{\infty}_{c}([-1,1])\) and set
\(\varphi_{A}(u)=A^{-1/2}\varphi(u/A)\). Then
\(\widehat{\varphi}_{A}(t)=A^{1/2}\widehat{\varphi}(At)\), so
\(\alpha(\varphi_{A})\to1\) as \(A\to\infty\). Since
\(\widehat{\varphi}\) is Schwartz,
\[
\sup_{T_{1}\geq t_{L}^{\ast}}
g_{L}(T_{1})\beta(\varphi_{A},T_{1})\to0 .
\]
For \(A\) large, \(g_{L}(T_{1})\beta(\varphi_{A},T_{1})<M\alpha(\varphi_{A})\)
for every \(T_{1}\). Thus \(\varphi_{A}\) lies outside every IE class.
The same concentration gives \(B[\varphi_{A}]\to0\) while
\(\sup_{|t|<t_{\ast}}|\widehat{\varphi}_{A}(t)|^{2}\to\infty\), so the
HFD condition also fails.
After multiplying by the standard even holomorphic cutoff used in
\Cref{v4-arith:w5-a:lem:upsilon-cuspidal}, the H1--H3 conditions are
preserved and the weighted-tail inequalities persist by openness.
\end{proof}

\begin{remark}[what the gap captures]
\label{v4-arith:w7-e:rem:gap-captures}
The gap
\(\mathcal{S}\setminus\bigcup_{T_{1}}\mathrm{PW}^{\mathrm{hol},\mathrm{IE}}_{F}(T_{1})\)
is non-empty. Its model elements are low-frequency approximate
identities with negligible logarithmic far tail. No assertion about a
particular Li approximant follows without computing its weighted
spectral profile.
\end{remark}

\section{Compatibility with Chapters 61, 57, and the Bounded-Interval Reduction}
\label{v4-arith:w7-e:compatibility}

\subsection{IE Class and Chapter 61 Bounded-Interval Reduction}
\label{v4-arith:w7-e:compat-ch61}

\begin{proposition}[IE class gives only the archimedean half of the bounded reduction]
\label{v4-arith:w7-e:prop:compat-bounded}
\ClaimStatusProvedHere. For every \(f\in\mathrm{PW}^{\mathrm{hol},\mathrm{IE}}_{F}(T_{1})\),
one has \(A[f]\leq B[f]\), hence \(W_{\infty}(f*\widetilde f)\ge0\).
The bounded-interval domination statement \((\star)\) of
\Cref{v4-arith:w6-a:thm:bounded-reduction} follows only after the
additional residual inequality
\(R[f]\geq W_{\infty}(f*\widetilde f)\).
\end{proposition}

\begin{proof}
\Cref{v4-arith:w7-e:thm:IE-ATP} gives \(A[f]\leq B[f]\). Substitution
into \Cref{v4-arith:w6-a:thm:bounded-reduction} shows that the remaining
condition is exactly \(B[f]\leq A[f]+2\pi R[f]\), equivalently
\(R[f]\geq W_{\infty}(f*\widetilde f)\).
\end{proof}

\subsection{Three Non-Elementary Constructions Revisited}
\label{v4-arith:w7-e:three-constructions}

The three non-elementary constructions of
\Cref{v4-arith:w6-a:rem:live-constructions} (zero-density refinement,
Gaussian reduction, Hermite--Biehler/de Branges) remain independent.
The IE theorem does not obviate them.

\begin{remark}[construction interactions]
\label{v4-arith:w7-e:rem:construction-interactions}
\begin{enumerate}
\item \emph{Zero-density refinement.} The IE class handles tail-heavy
PW; the Riemann--von Mangoldt asymptotic density
\(N(T)\sim(T/2\pi)\log T\) controls the density of zero ordinates, but
it does not turn the zero functional into a positive form. For
IE-captured \(f\), A-TP is obtained before appealing to a zero-side
interpretation. For residual \(f\), the obstruction is precisely the
Weil positivity problem; it cannot be removed by calling the
zero-side term positive.
\item \emph{Gaussian reduction.} Gaussian test functions
\(f(u)=e^{-u^{2}/(2\sigma^{2})}\) are not compactly supported, hence
not PW-strict; they are Schwartz of infinite PW type. Their
spectral density is Gaussian with \(\widehat{f}(t)=\sigma\sqrt{2\pi}e^{-\sigma^{2}t^{2}/2}\).
For \(\sigma\ll 1\), spectrally wide: \(\alpha\) small, \(\beta\) large,
IE condition holds. For \(\sigma\gg 1\), spectrally narrow:
\(\alpha\) close to 1, IE fails. The Hardy--Littlewood Gaussian-Weil
obstruction pursues the narrow-Gaussian regime, a model for the IE
residual. See
Chapter~\ref{v4-arith:w5-i:chapter}.
\item \emph{de Branges structure.} The prolate-eigenvalue
characterization \Cref{v4-arith:w7-e:thm:prolate-characterization}
is a time-frequency estimate in the Paley--Wiener model. The de
Branges construction asks for positivity in \(\mathcal{H}(E_{\xi})\). The two
structures are compatible, but the prolate estimate is not a
substitute for Hermite--Biehler positivity.
\end{enumerate}
The unconditional \(W_{\infty}\)-positivity on
\(\mathrm{PW}^{\mathrm{hol},\mathrm{IE}}_{F}\) adds a far-tail
sufficient criterion. Full A-TP still retains the \(R[f]\) term. The
two-scale prolate inequality gives a separate uniform sufficient
criterion when its numerical inequality is verified.
\end{remark}

\section{Bounded-Direct Boundary}
\label{v4-arith:w7-e:summary}

\begin{enumerate}
\item The digamma kernel \(g(t)=\mathrm{Re}\,\psi(1/4+it/2)+\log\pi\)
admits an explicit Stirling lower bound
\(g(t)\geq g_{L}(t)=\tfrac{1}{2}\log((1+4t^{2})/16)+\log\pi-2/(1+4t^{2})\)
for \(|t|\geq t_{L}^{\ast}\)
(\Cref{v4-arith:w7-e:thm:g-lower-bound}), with \(g_{L}\) vanishing at
\(t_{L}^{\ast}\approx 0.8946>t_{\ast}\approx 0.8507\)
(\Cref{v4-arith:w7-e:cor:g-L-zero}).
\item The integral-energy sub-class
\(\mathrm{PW}^{\mathrm{hol},\mathrm{IE}}_{F}(T_{1})\) is defined by
\(M\alpha(f)\leq g_{L}(T_{1})\beta(f,T_{1})\)
(\Cref{v4-arith:w7-e:def:IE-subclass}), with \(\alpha,\beta\) the
core and far-tail \(L^{2}\)-energy fractions.
\item The archimedean term \(W_{\infty}\) is non-negative on
\(\mathrm{PW}^{\mathrm{hol},\mathrm{IE}}_{F}(T_{1})\) for every
\(T_{1}\geq t_{L}^{\ast}\)
(\Cref{v4-arith:w7-e:thm:IE-ATP}); full A-TP also requires
\(R[f]\geq W_{\infty}(f*\widetilde f)\).
\item The union of IE with the bounded-interval reduction HFD class is a strict
enlargement of HFD: IE contains narrow-core, distributed-tail profiles
which HFD misses, but HFD is not contained in IE
(\Cref{v4-arith:w7-e:prop:strict-inclusion}).
\item The prolate-eigenvalue bound
\(\alpha(f)\leq\lambda_{0}(T,t_{\ast})\) must be paired with the
far-tail bound
\(\beta(f,T_{1})\geq1-\lambda_{0}(T,T_{1})\). The resulting two-scale
inequality gives a uniform \(W_{\infty}\)-positivity result only when
\eqref{v4-arith:w7-e:eq:two-scale-prolate-condition} is verified
(\Cref{v4-arith:w7-e:cor:uniform-narrow-PW}).
\item The IE residual is the locus where
\(g_{L}(T_{1})\beta(f,T_{1})<M\alpha(f)\) for every \(T_{1}\)
(\Cref{v4-arith:w7-e:thm:IE-residual}). It is a weighted far-tail
failure, not a claim about a named Li approximant.
\item Strict inclusions:
\[
\mathrm{PW}^{\mathrm{hol},\mathrm{HFD}}_{F}\subsetneq
\mathrm{PW}^{\mathrm{hol},\mathrm{HFD}}_{F}\cup
\bigcup_{T_{1}}\mathrm{PW}^{\mathrm{hol},\mathrm{IE}}_{F}(T_{1})
\subsetneq\mathrm{PW}^{\mathrm{hol}}_{F}.
\]
(\Cref{v4-arith:w7-e:thm:extent}).
\end{enumerate}

\begin{remark}[far-tail Beilinson principle]
\label{v4-arith:w7-e:rem:beilinson-closing}
The integral-energy reduction obeys the Beilinson principle: a
transverse sufficient criterion is an honest substitute for no
inclusion theorem. The estimate is explicit, elementary, and
computable: the constants
\(M\approx 3.083\), \(t_{\ast}\approx 0.8507\),
\(t_{L}^{\ast}\approx 0.8946\), and the Stirling lower bound \(g_{L}\)
are all closed-form. The residual is precisely the failure of all
logarithmic far-tail payments.
\end{remark}

\begin{remark}[quantification of the bounded step]
\label{v4-arith:w7-e:rem:quantification}
Relative to the bounded-interval reduction HFD (sup-bound sub-class), the integral-energy reduction proves
\(W_{\infty}\ge0\) on:
\begin{itemize}
\item All \(f\) satisfying a verified two-scale prolate inequality
(\Cref{v4-arith:w7-e:cor:uniform-narrow-PW}).
\item All \(f\) satisfying the integral-energy condition
\(M\alpha\leq g_{L}(T_{1})\beta\) for some \(T_{1}\).
\end{itemize}
The uncovered residual is the failure of the weighted far-tail
inequality. Bridging this residual requires a separate positivity
mechanism, such as a proved Gaussian reduction or the de Branges
Hermite--Biehler construction. The elementary ceiling reached by
integral-energy bounds is the IE sub-class.
\end{remark}
